\paragraph{Order \(\Lambda^2\): contacts, vector holonomy and the Noether coefficient of \(\rho^2X\) against the gap of the normalisation paragraph (conditional; not closure).}
Keep the setting of the paragraph on the normalisation of the tensor kinetic term and of the paragraphs on the logarithmic coefficient and its contact terms: Poisson limit, Euclidean section, \(\alpha=1\), one loop, the \(\LUV^2\) term, a totally geodesic \(\Sigma\), cross terms at first order in \(1-n\), the family \(\Theta^{12}=1/f(x^1,x^2)\), \(\Theta^{34}=1\), \(f=\lvert z\rvert^{n-1}\), and \(\sigma=+1\) for \(n<1\), \(\sigma=-1\) for \(n>1\). Write the entropy of a sector as \(S=(\LUV^2A_G/8\pi)\,c\), so that a flat-cone scalar component has \(c=\tfrac16\) and \(1/G=c\,\LUV^2/2\pi\). Then \(c_A=N_{\mathrm{tot}}/6\) (the antiperiodic Dirac gluing adds \(n_D\)) and \(c_T=N_{\mathrm{tot}}/6-4\kappa_X\), the values of the paragraphs on the normalisation and on fermionic content. Before any contact, \(c_A-c_T=4-n_D\) for the fixed-state replica and \(4-2n_D\) for the periodic one; at \(n_D=0\) this is \(1/G_A-1/G_T=2\Lambda_{\mathrm{UV},v}^2/\pi\), nonzero and independent of \(N_s\).

\emph{Contacts.} The coefficient \(a_1=\operatorname{tr}(R/6-E)\) is linear in the endomorphism, so its cross term with the deltas is read off from the delta content of the family. The uniform \(E_\delta\) of the contact paragraph lies on every scalar, connection, ghost and Dirac component and shifts \(c\) by \(-N_{\mathrm{tot}}/4\); it carries the factor \(N_s\), and with the scalars at \(\Lambda_{\mathrm{UV},s}\) and the connection with its ghosts at \(\Lambda_{\mathrm{UV},v}\) the residue is \(\tfrac72\Lambda_{\mathrm{UV},v}^2-\tfrac14N_s\Lambda_{\mathrm{UV},s}^2\) at \(n_D=0\), which cannot follow a gap that does not depend on \(N_s\). The holonomy \(\alpha_v\mathcal J\) of the vector bundle around \(\Sigma\) twists the two components with \(\mathcal J=\pm1\) by \(a=\alpha_v\lvert1-n\rvert\). The constant term of the heat trace of one component on the cone of angle \(2\pi n\) with phase \(e^{2\pi\mathrm ia}\) is \(C(n,a)=B_2(a)/(2n)-n/12\), \(B_2(a)=a^2-a+\tfrac16\), \(0\le a\le1\); it gives \((1/12)(1/n-n)\) at \(a=0\) and the Aharonov--Bohm value \(-a(1-a)/2\) at \(n=1\), and agrees to \(10^{-12}\) with a numerical sum over the Bessel orders \(\lvert m+a\rvert/n\) at twelve pairs \((n,a)\). Its term linear in \(a\) shifts \(c\) of the vector by \(-\alpha_v\sigma\): \(-1\) for \(\alpha_v=1\) on \(n<1\), the value that takes the vector with its ghosts from \(\tfrac13\) to \(-\tfrac23\); the cone angle of the family fixes \(\alpha_v=1\). For a Dirac field with the less singular extension the constant is \(a^2+\tfrac16\) per block and the symmetric twists \(a_k=k/n\) sum to zero, so a common holonomy shift has no term of first order: the Dirac holonomy contact is zero at this order. At \(n_D=0\) and \(\alpha_v=1\) the residue \(c-c_T\) is
\begin{center}
\begin{tabular}{@{}lc@{}}
\toprule
replica and contacts & \(c-c_T\)\\
\midrule
fixed-state, no delta, no holonomy & \(4\)\\
fixed-state and vector holonomy, \(n<1\) / \(n>1\) & \(3\) / \(5\)\\
uniform \(E_\delta\) only & \(\tfrac72-\tfrac14N_s\)\\
\(E_\delta\) and holonomy, \(n<1\) / \(n>1\) & \(\tfrac52-\tfrac14N_s\) / \(\tfrac92-\tfrac14N_s\)\\
\bottomrule
\end{tabular}
\end{center}
with the single cutoff of the normalisation paragraph. The first two rows do not vanish for any \(N_s\); the last two vanish at \(N_s=14\), \(10\), \(18\) only, a choice of content. For \(n_D>0\) the first two rows become \(4-n_D\) and \(4-n_D-\alpha_v\sigma\), which vanish at \(n_D=4\) and \(n_D=3\) (\(n<1\)) for every \(N_s\), again a choice of content, and equal \(2\) and \(1\) at \(n_D=2\). A holonomy with \(\alpha_v\sigma=4\) would close the bosonic gap; it needs a second circulation of modulus \(3\) in \(\Theta^{34}\), which makes the metric induced on \(\Sigma\) singular, and it is not used. The statement of the normalisation paragraph that \(\rho^2X\) contributes no entropy refers to the fixed-state replica, with \(\rho\) regular, or equivalently to the presentation \(r=0\) below; a Wald-type coefficient of \(\rho^2X\) is presentation dependent.

\emph{Dragging, regularity, covariance.} Take \(e^A{}_\mu=\eta^{AB}\omega_{\mu B}\) as in the paragraph on the torsion invariant. Let \(\xi\) be the rotation (Lorentzian boost) of the normal plane with generator \(\mathcal J\). The family, for any \(f(r)\), and every constant \(\bar\Theta\) with \([\mathcal J,\bar\Theta]=0\) (the span of \(e_{12}\) and \(e_{34}\), two of the six components) are invariant under the flow that rotates the chart and the matrix index together, and not under the rotation of the chart alone: for \(\bar\Theta=a(e_{12}+e_{34})\) the Lie derivative along the chart rotation is \(\pm a\,\mathrm{diag}(1,1,0,0)\), the sign following the orientation of \(\mathcal J\) (exact, both signatures). The background \(n=1\) is regular; for \(n\ne1\), \(\Theta^{12}=1-(n-1)\ln r+O((n-1)^2)\) and \(X=2(n-1)^2r^{-2n}\) are singular on \(\Sigma\); this singular-\(\rho\) family is the second reading of the paragraph on the geometry that carries the cone. The Noether-charge construction of~\cite{IyerWald1994} needs a diffeomorphism covariant Lagrangian of the dynamical fields, stationary fields with \(\mathcal L_\xi\phi=0\), and regularity on the bifurcation surface; the term \(Y(\phi,\mathcal L_\xi\phi)\) by which the charge is ambiguous is linear in \(\mathcal L_\xi\phi\) and drops out for stationary fields, which is the step that makes the entropy independent of the ambiguities~\cite{JacobsonKangMyers1994}. Here \(\mathcal L_\xi e^A{}_\nu=e^A{}_\sigma\partial_\nu\xi^\sigma\) on \(\Sigma\), so \(e_{A\mu}\mathcal L_\xi e^A{}_\nu=\nabla_\nu\xi_\mu\ne0\): the frame is dragged only up to a rotation of the matrix index, a global symmetry of the induced density. This is the form of the obstruction for the unit vector of Einstein-aether theory, where the derivation fails and no entropy has been identified~\cite{Foster2006}; here a compensating rotation exists but is rigid. The contraction of the matrix index is with the fixed \(\eta\) of the action, not with a spacetime metric.

\emph{Noether coefficient.} In the coframe variables the density \(\tfrac{N_{\mathrm{tot}}}6\rho R[g_\rho]+\kappa_X\rho^2X\), with \(X=T_{\lambda\mu\nu}T^{\lambda\mu\nu}\), is a diffeomorphism covariant first-order functional of \(e^A{}_\mu\) (through \(\det e\) and \(\mathrm de\)). For every such functional \(J^\mu=\partial_\nu Q^{\mu\nu}-E^\mu{}_Ae^A{}_\sigma\xi^\sigma\) with \(Q^{\mu\nu}=(\partial L/\partial\partial_\mu e^A{}_\nu)\,e^A{}_\sigma\xi^\sigma\) (checked at three random fields to \(10^{-30}\)), which vanishes where \(\xi=0\): the torsion invariants carry no Noether charge on \(\Sigma\), and the Noether coefficient is the coefficient of \(R\) in the presentation chosen, where \(R\) denotes \(R[G]\) with \(\sqrt G=\rho^2\); the presentations differ from \(\rho R[g_\rho]\) by terms in \(V^2\) and total derivatives, which are kept in the remainder. It is not unique, because \(R=-\mathcal T-2\nabla_\mu V^\mu\) with \(\mathcal T=\tfrac14T_{\lambda\mu\nu}T^{\lambda\mu\nu}+\tfrac12T_{\lambda\mu\nu}T^{\mu\lambda\nu}-V_\mu V^\mu\) and \(V_\mu=\partial_\mu\ln\rho\) (exact at three random two-jets, with the identity of the paragraph on the torsion invariant). With the vector, axial and tensor parts \(V,S,q\) of the torsion, \(X=\tfrac23V^2+S^2+q^2\) and \(\mathcal T=-\tfrac23V^2-\tfrac14S^2+\tfrac12q^2\). Eliminating \(q^2\), \(V^2\) or \(S^2\) gives the coefficient \(r=-2\), \(1\), \(4\) of \(R\) in \(X\), and the variables \((G,\rho,b)\) give \(r=0\); any real \(r\) is obtained by adding \(r(R+\mathcal T+2\nabla_\mu V^\mu)=0\). The Noether coefficient is \(c_W(r)=N_{\mathrm{tot}}/6+\kappa_Xr\), so \(c_W-c_T=\kappa_X(r+4)\), zero for every content only at \(r=-4\). On the sampled modes with \(\delta\rho=0\) of both chiralities, \(S^2:q^2:X=\tfrac13:\tfrac23:1\) and \(\langle R\rangle=-\tfrac14\langle X\rangle\); the presentation \(r=-4\) is the one whose torsion remainder, \(\tfrac{10}3V^2+2S^2-q^2\), vanishes on every such mode and not on pairs of modes with different wave vectors. The value \(r=-2\) of the paragraph on the torsion invariant leaves \(c_W-c_T=2\kappa_X\), the part of the on-mode value of \(X\) that its remainder \(2V^2+\tfrac32S^2\) carries. Among \(aT_{\lambda\mu\nu}T^{\lambda\mu\nu}+bT_{\lambda\mu\nu}T^{\mu\lambda\nu}+cV^2\) the antisymmetric part of the frame equation vanishes identically only for \((a,b,c)\propto(\tfrac14,\tfrac12,-1)\) (numerical, \(d=3\), residual \(10^{-29}\)), so for \(\kappa_X\ne0\) the density is not invariant under local rotations of the frame, and the class of \(X\) modulo that direction is canonical.

\emph{Cross term with the connection of the cone.} On the family the connection of the vector bundle is \(b=(1-n)\mathcal J\,\mathrm d\theta\), with holonomy \(2\pi(1-n)\mathcal J\) and no delta. The cross term of \(\rho^2X\) with a regular first-order connection \(\beta\) is \((n-1)\int\Omega_{12}\ln(L/r)\,\mathrm d^2x\) up to normalisation (exact for polynomial \(\beta\)): the weight is the Green function of the normal plane and not \(\delta_\Sigma\), so it is a bulk term of the kind removed by the equations of motion in the construction of~\cite{LewkowyczMaldacena2013}, not an area term. The charge of the constraint term of the composite connection is not needed for this dependence, which is algebraic in \(b\), but it is not computed.

None of the contacts above (E-type, \(-\alpha_v\sigma\) with \(\alpha_v=1\), zero for Dirac) supplies the required \(-4\kappa_X\) of \(c\): the available vector contact has modulus \(1\) and the \(E_\delta\) contact depends on \(N_s\). The Noether coefficient is not obstructed but undetermined: the density fixes no value of \(r\), the value \(r=-4\) would match and is selected only by the match, and the value \(r=-2\) of the paragraph on the torsion invariant leaves \(2\kappa_X\). It stays undetermined until the frame is dragged by a flow whose internal part is a local symmetry, that is until the question of local frame rotations (Problem~D) is settled. These statements are for the Poisson-limit operators, one loop, \(\alpha=1\), the \(\LUV^2\) term, the family and the first-order additive cross-term rule on the \(n<1\) and \(n>1\) branches; they do not derive the replica, the holonomy or the state selected by the matrix model, the first law with the rigid compensating rotation, the Dirac extension, the exact spectrum of the singular cone with holonomy, or a value or sign of \(G_{\mathrm{eff}}\). Problem~B stays open and PROGRAMME-CONDITIONAL.
